\documentclass{article}
\usepackage{graphicx}
\usepackage{amsmath}
\usepackage{hyperref}

\title{Failure Propagation in Multi-Agent LLM Systems: Characterization and Containment with Circuit Breakers}
\author{Goutham Patley \\ Independent Researcher \\ gp4em@virginia.edu}
\date{}

\begin{document}

\maketitle

\begin{abstract}
Large language models (LLMs) are increasingly deployed as compositional systems in which multiple agents interact to solve complex tasks. While such architectures enable modularity and improved capability, they introduce a new class of failure modes: errors originating in one component can propagate through the system, leading to cascade failures.

We study failure propagation in multi-agent LLM pipelines and evaluate circuit breakers as a mechanism for containment. We introduce Cascade Failure Rate (CFR) and provide both simulation and real-world validation.

Our results show that up to 62\% of upstream failures propagate to final failure without protection, while circuit breakers can eliminate cascade failures entirely. We further show that adaptive containment strategies introduce a tradeoff between failure containment and task success.

These findings suggest that reliability in compositional LLM systems is governed not only by component-level accuracy but by the dynamics of failure propagation across stages.
\end{abstract}

\section{Introduction}

Multi-agent LLM systems are increasingly used for planning, tool use, and reasoning. These systems improve capability but introduce new failure modes.

A key challenge is failure propagation: errors introduced in early stages can cascade through downstream components. Unlike isolated errors, these failures compound across stages.

Failure propagation scales multiplicatively with pipeline depth, making system reliability fundamentally different from component-level reliability.

We frame compositional LLM systems as distributed systems where reliability depends on controlling error propagation. Inspired by circuit breakers~\cite{hystrix,nygard}, we study mechanisms to contain failures before they cascade.

\textbf{Contributions}
\begin{itemize}
\item Formalize failure propagation using Cascade Failure Rate (CFR)
\item Introduce circuit breaker mechanisms for LLM pipelines
\item Provide simulation and real-world validation
\item Analyze containment vs success tradeoffs
\end{itemize}

\section{Theoretical Framework}

Let $p_u$ denote upstream failure probability and $p_p$ denote propagation probability.

\begin{equation}
P(\text{final failure}) = p_u \cdot p_p
\end{equation}

\begin{equation}
\text{CFR} = P(\text{final failure} \mid \text{upstream failure})
\end{equation}

\textbf{Proposition.} For a pipeline of length $L$:
\begin{equation}
\text{CFR}(L) = 1 - (1 - p_p)^L
\end{equation}

which is monotonically increasing in $L$.

\textbf{Implication.} Even moderate propagation probabilities lead to high end-to-end failure rates in deeper pipelines. Circuit breakers reduce effective $p_p$ by interrupting propagation.

\section{Method}

We evaluate three configurations:
\begin{itemize}
\item No Circuit Breaker
\item AI-based Circuit Breaker
\item Adaptive Circuit Breaker
\end{itemize}

\textbf{Pipeline Configuration.} We use a 3-stage pipeline consisting of a planner, executor, and verifier. Each stage produces structured outputs consumed by downstream stages.

\textbf{Failure Injection.} Failures are injected according to upstream failure probability $p_u$, while propagation behavior is governed by $p_p$.

\textbf{Failure Definition.} An upstream failure is defined as:
\begin{itemize}
\item invalid or malformed output
\item missing required fields
\item semantically inconsistent response
\end{itemize}

\textbf{Circuit Breaker Policy.} The AI-based circuit breaker evaluates intermediate outputs using a secondary model call. The adaptive circuit breaker dynamically adjusts its triggering sensitivity based on recent failure signals, increasing intervention probability following detected failures.

\textbf{Metrics.} We report:
\begin{itemize}
\item Cascade Rate (fraction of upstream failures that propagate)
\item Final Success Rate (fraction of tasks completed successfully)
\item Breaker Activation Rate (average number of interventions)
\end{itemize}

\section{Experiments}

\subsection{Simulation}

We simulate pipelines of varying depth (2--6 stages). Each stage introduces failures with probability $p_u$, and failures propagate with probability $p_p$.

\subsection{Real-world API Setup}

We evaluate a planner–executor–verifier pipeline using a GPT-4-class LLM API. Each experiment consists of 50 independent runs on a synthetic task set requiring structured reasoning and multi-step execution.

Tasks involve generating intermediate structured outputs that are consumed by downstream stages. Failures are injected probabilistically at intermediate stages, and propagation is observed through to final task outcomes.

An upstream failure is defined as an intermediate output that is invalid, incomplete, or low-confidence. A cascade occurs when such a failure leads to final task failure.

\section{Results}

\subsection{Simulation Results}

\begin{figure}[h]
\centering
\includegraphics[width=0.7\linewidth]{figure_simulation_cfr.png}
\caption{CFR increases with pipeline depth}
\end{figure}

CFR increases monotonically with pipeline depth, consistent with the theoretical model. This empirically validates that failure propagation grows with system depth.

Circuit breakers significantly reduce propagation, flattening the growth of CFR and effectively lowering the system-level propagation probability.

\subsection{Real-world Results}

All results are averaged over 50 runs. Variance across runs was low (standard deviation $< 0.05$), indicating stable behavior across configurations.

\begin{table}[h]
\centering
\begin{tabular}{lcc}
\hline
Method & Cascade Rate & Success \\
\hline
No CB & 0.62 & 0.38 \\
AI CB & 0.00 & 0.44 \\
Adaptive CB & 0.00 & 0.40 \\
\hline
\end{tabular}
\caption{Real-world results}
\end{table}

\begin{figure}[h]
\centering
\includegraphics[width=0.7\linewidth]{figure_real_api_results.png}
\caption{Real-world validation}
\end{figure}

Without circuit breakers, 62\% of upstream failures propagate to final failure, demonstrating significant cascade effects in real systems.

Both AI-based and adaptive circuit breakers eliminate cascade failures entirely. However, the AI-based circuit breaker achieves higher final success rates, indicating that selective intervention is more effective than aggressive suppression.

\section{Tradeoff Analysis}

\begin{figure}[h]
\centering
\includegraphics[width=0.6\linewidth]{figure_tradeoff.png}
\caption{Containment vs success tradeoff}
\end{figure}

Adaptive policies eliminate cascades but reduce success rates due to over-intervention. This occurs because they prioritize containment without jointly optimizing for task completion.

This tradeoff suggests that optimal reliability mechanisms must balance failure sensitivity with tolerance for recoverable intermediate errors. Excessively aggressive intervention suppresses useful outputs and reduces overall system effectiveness.

This behavior is analogous to over-triggering in traditional circuit breakers, where excessive sensitivity reduces throughput. Effective policies must therefore jointly optimize containment and successful task execution.

\section{Discussion}

Failure propagation is a system-level phenomenon. Reliability depends on controlling propagation dynamics rather than only improving individual components.

Our results suggest that compositional LLM systems should be designed with explicit failure containment mechanisms. This shifts the focus from improving individual model outputs to managing system-level error dynamics.

The behavior of adaptive strategies further highlights that containment alone is insufficient. Minimizing cascade failures without accounting for task success leads to excessive intervention and degraded performance.

\section{Related Work}

Multi-agent systems such as AutoGPT~\cite{autogpt} and BabyAGI~\cite{babyagi} demonstrate compositional LLM architectures. Toolformer~\cite{toolformer} and ReAct~\cite{react} explore tool use and reasoning.

Error propagation in machine learning pipelines has been studied in prior work~\cite{sculley2015hidden}, which highlights how small upstream issues can amplify downstream.

Recent work on LLM reliability has focused on reasoning robustness and self-consistency~\cite{wang2022selfconsistency}.

Distributed systems literature~\cite{hystrix,nygard} motivates circuit breaker design.

\section{Limitations}

We evaluate relatively simple pipelines and controlled failure injections. Real-world systems may exhibit more complex dependencies.

\section{Future Work}

Future work should explore principled adaptive policies and extend to larger multi-agent systems.

\bibliographystyle{plain}
\bibliography{references}

\end{document}
